\chapter{Capstone: One Company, Every Lens}\label{ch:capstone}

% Scope: the running B2B SaaS company walked through every Part of the book in sequence,
%   tracing how each discipline feeds the next along the value-creation spine.
%   No new eq: labels. Cross-references only. One figure: value-spine.
%   Structure: framing paragraph → fact-sheet table → 9 sections (one per Part + new 2026 additions).

Every framework in this book was introduced independently. Here they converge. The same
\money{4000000}-ARR B2B SaaS company that appeared in every worked example is now walked
end-to-end across all 9 Parts and 32 chapters: from the first test of value creation and the
accounting substrate that grounds every number, through customer economics and paid acquisition,
into competitive position and strategic durability, then the financial model that prices the
whole, and finally the execution disciplines that determine whether the plan survives contact
with a quarter. The purpose is not repetition --- each section assumes its chapter has been
read --- but integration: the connections the separate chapters could only promise are now
explicit. The seven operational Parts (Foundations, Customer, Acquisition, Strategy, Finance,
Execution, and the Founder) are followed by a Reference Part and this Capstone; the eleven
chapters added in the 2026 edition --- covering financial statements, experiments, product
management, channels and go-to-market, data and AI strategy, the operating model, fundraising
mechanics, and the founder's own craft of persuasion, storytelling, psychology, and leadership
--- are woven into every section below.

\medskip
\noindent\textbf{Running-company fact sheet.}

\medskip
\noindent\begin{tabularx}{\linewidth}{@{}l >{\raggedright\arraybackslash}X@{}}
  \toprule
  Parameter & Value \\
  \midrule
  Plan price (ARPU)                 & \money{50}/mo (\money{600}/yr) \\
  LTV (growing perpetuity)          & \money{3150} \\
  Contribution margin               & \qty{42}{\percent} → \money{21}/mo \\
  Target CAC (operating)            & \money{600} → LTV:CAC 5.25× \\
  Steady-state growth $g$           & \qty{3}{\percent}/yr \\
  Monthly ad budget → new customers & \money{120000} → 200/mo \\
  WACC / discount rate              & \qty{11}{\percent}/yr \\
  Funnel                            & 10,000 sessions → 800 trials → 200 paid \\
  NOPAT (approx.)                   & \money{600000} \\
  ARR / MRR                         & \money{4000000} / \money{340000} \\
  Invested capital                  & \money{3000000} \\
  ROIC                              & \qty{20}{\percent} \\
  Beta ($\beta$)                    & 1.4 → cost of equity \qty{11}{\percent} \\
  Gross margin                      & \qty{70}{\percent} \\
  Year-3 EBITDA margin              & \qty{+7}{\percent} (Rule-of-40 score: 37) \\
  Series-A pre-money (market)       & \money{24000000} (\(6\times\) ARR × Rule-of-40 adj.) \\
  Series-A pre-money (DCF floor)    & \money{15000000} (base scenario, \cref{eq:scenario-ev}) \\
  Runway at current burn            & 8 months (\money{2500000} cash ÷ \money{310000} net burn) \\
  SAFE outstanding                  & \money{500000} at \money{5000000} cap, \qty{20}{\percent} discount \\
  \bottomrule
\end{tabularx}
\medskip

% complexity: 5 -> figure value-spine

\begin{figure}[htbp]
  \centering
  \includegraphics[width=0.96\linewidth]{value-spine}
  \caption{The value-creation spine. Each stage of the book feeds the next: Foundations set
    the ROIC\,>\,WACC test; Customer economics set the CLV that caps acquisition spend; Acquisition
    validates channel efficiency via ROAS and iROAS; Strategy sets the moat-driven growth rate $g$
    used in the terminal value; Finance discounts the resulting cash flows at WACC; Execution
    disciplines (Rule of 40, quick ratio) confirm the plan is being realised; the Founder skills
    (influence, story, leadership) turn the plan into closed customers, a funded round, and a led
    team; the Reference chapter closes the measurement loop.}
  \label{fig:value-spine}
\end{figure}

%% ─────────────────────────────────────────────────────────────
\section{Foundation: Does It Create Value?}\label{sec:capstone-value}
%% ─────────────────────────────────────────────────────────────

The master test is \cref{eq:economic-profit}: economic profit equals NOPAT minus the capital
charge. The \money{600000} NOPAT, \money{3000000} invested capital, and \qty{20}{\percent}
ROIC are not assumptions --- they are derived from the three-statement model of
\cref{ch:financial-statements}. \Cref{sec:three-statements} and \cref{sec:statement-analysis}
trace every dollar from the income statement through the NOPAT bridge of \cref{eq:nopat-bridge}:
EBIT \(\times (1-30\%) = \money{860000}\times0.70 \approx \money{600000}\). Operators who
cannot trace those numbers back to accounting statements are working from faith. The inference
discipline of \cref{ch:experiments-and-inference} underlies every causal claim in the book:
the statistical machinery of \cref{eq:z-test} and \cref{eq:sample-size} determines whether
an observed ROIC improvement is real or sampling noise, and the OLS framework of \cref{eq:ols}
controls for the covariates that would otherwise confound attribution.

With NOPAT \(\approx\money{600000}\) and invested capital of \money{3000000},
\cref{eq:roic} gives \(\mathrm{ROIC}=20\%\). WACC, derived in \cref{ch:risk-and-capital-structure}
via \cref{eq:wacc} and \cref{eq:capm} at \(\beta=1.4\), is \qty{11}{\percent}:
\[
  \mathrm{EP}
    = (0.20 - 0.11)\times\money{3000000}
    = \qty{9}{\percent}\times\money{3000000}
    = \money{270000}.
\]
The spread is positive and material. A \qty{9}{\percent} ROIC--WACC gap is not a rounding
error; it is evidence that the business earns \money{270000} per year above the investors'
opportunity cost. Every dollar of additional invested capital deployed at the current ROIC
creates \money{0.09} of annual economic profit --- a compounding claim, not a one-time gain.

The implication of \cref{eq:economic-profit} runs in both directions. If WACC rises --- say,
because rising rates push \cref{eq:capm} from \qty{11}{\percent} to \qty{13}{\percent} ---
the spread compresses from \qty{9}{\percent} to \qty{7}{\percent} and EP falls to \money{210000}
without any operational deterioration. Conversely, cutting invested capital by \money{500000}
(through tighter working-capital management) while holding NOPAT constant raises ROIC to
\qty{24}{\percent} and EP to \money{325000}. The capital denominator is as important as the
numerator.

This verdict from \cref{ch:value-creation} frames every subsequent analysis. The chapters
that follow each translate this single question into their own terms: CLV:CAC is its
per-customer expression (\cref{ch:customer-economics}); the DCF enterprise value is its
cumulative present value (\cref{ch:valuation}); the moat is the mechanism that sustains it
(\cref{ch:advantage-and-disruption}). \textcite{brealey2020} and \textcite{damodaran2012}
develop the corporate-finance architecture behind \cref{eq:economic-profit,eq:roic,eq:wacc}.

%% ─────────────────────────────────────────────────────────────
\section{Customer Understanding: Who, and How Much?}\label{sec:capstone-customer}
%% ─────────────────────────────────────────────────────────────

Value creation begins with a choice of whom to serve. \cref{eq:segment-score} applied in
\cref{ch:segmentation-targeting-positioning} scores mid-market (50--500 seats) at 7.70
versus SMB at 6.20 across market size, competitive intensity, LTV potential, and sales-cycle
fit. The decisive criterion is LTV potential: mid-market ARR is roughly \(10\times\) the
SMB level, which cascades into every downstream number.

With the target segment fixed, the pricing headroom is bounded by the economic value to
customer (\cref{eq:evc} in \cref{ch:pricing}): the product's contribution to the
customer's outcome sets the theoretical price ceiling above which willingness to pay
collapses. Our \money{50}/month price sits well below that ceiling for mid-market buyers ---
the primary justification for the product-development move toward a \money{299}/month enterprise
tier explored in \cref{ch:growth}.

The per-customer financial case is \cref{eq:clv}. At \qty{42}{\percent} contribution margin
(\money{252}/year), WACC \qty{11}{\percent}, and steady-state margin growth \qty{3}{\percent}:
\[
  \mathrm{CLV} \;=\; \frac{\money{252}}{0.11 - 0.03} \;=\; \money{3150}.
\]
This is the ceiling on rational acquisition spend. \cref{eq:ltv-cac} converts it into the
operating constraint: at a \money{600} target CAC, LTV:CAC\,=\,5.25×, well above the
3× floor that proxies the ROIC\,>\,WACC test at the customer level. The gap between
5.25× and the 3× floor represents the buffer that absorbs model uncertainty in CLV ---
the practical equivalent of the margin of safety in \cref{ch:corporate-finance-core}.

The brand lens from \cref{ch:brand-growth} adds a demand-side observation: brand salience
reduces the mental effort buyers expend at category entry points, lowering effective
CAC at the margin. A firm that earns the \money{3150} CLV but spends only \money{400} in
CAC because buyers arrive already primed has a structurally superior unit-economics position
--- the same ROIC test, run at lower capital cost.

\Cref{ch:product-management} establishes \emph{what} to build before pricing sets what to
charge. The enterprise tier --- the \money{299}/month product expansion referenced in
\cref{ch:pricing} --- did not emerge from intuition. It passed the RICE prioritisation gate
of \cref{eq:rice}, which scored it higher than alternatives by reaching more users with
higher confidence per unit of engineering effort. The Bass launch model (\cref{eq:bass})
then forecasts how quickly the enterprise segment will absorb the new tier: with an addressable
ceiling of 500 mid-market accounts and an imitation coefficient of \(q=0.3\), new-adopter
volume peaks around month 10, before which outreach must be front-loaded to seed the social
term. Product-market fit itself is measured through \cref{eq:retention}: the asymptotic
survival curve that confirms genuine fit is the same cohort-retention curve that drives CLV.

Together, \cref{ch:segmentation-targeting-positioning}, \cref{ch:brand-growth},
\cref{ch:pricing}, \cref{ch:product-management}, and \cref{ch:customer-economics} answer
the customer half of the value equation: who creates \money{3150} of value, what to build to
sustain it, and at what acquisition cost.
\textcite{gupta2004customer} is the authoritative treatment of CLV as a firm-value driver.

%% ─────────────────────────────────────────────────────────────
\section{Acquisition: Buying Demand Profitably}\label{sec:capstone-acquisition}
%% ─────────────────────────────────────────────────────────────

Customer economics sets what a conversion is worth; acquisition channels determine
what one costs. \cref{eq:funnel-chain} in \cref{ch:digital-funnel-and-crm} traces the
conversion cascade: 10,000 sessions produce 800 free trials at an \qty{8}{\percent}
trial rate, and 200 of those trials convert to paid at \qty{25}{\percent}. The overall
conversion rate is \qty{2}{\percent}. The binding constraint at steady state is the
trial-to-paid gate, not the top of funnel --- a theory-of-constraints
finding that \cref{ch:operations} generalises: elevating the non-binding session step
from \qty{8}{\percent} to \qty{10}{\percent} adds only 50 additional customers per
month; elevating the trial-to-paid rate from \qty{25}{\percent} to \qty{35}{\percent}
adds 80.

On the paid side, \cref{eq:roas-cpa} from \cref{ch:paid-acquisition-platforms} governs
the efficiency test. A target CPA of \money{600} at \money{120000}/month delivers 200
conversions --- precisely the learning-phase volume the platform delivery algorithms require
to stabilise. Setting the target at \money{300} to ``improve efficiency'' halves the
event volume, traps the ad set in Learning Limited, and causes realised CPA to rise rather
than fall: the platform cannot estimate \(\hat{p}\) accurately below roughly 50 conversion
events per week.

The harder question is whether the \money{600} CPA reflects truly incremental customers or
includes organic demand that would have converted anyway. \cref{eq:iroas} from
\cref{ch:attribution-and-measurement} answers it via geo-holdout or ghost-bidding
experiments: the incremental ROAS isolates the lift attributable to paid media, stripped
of last-click inflation. A business running at iROAS substantially below its
target-ROAS assumption is subsidising organic demand with paid budget and overstating
channel efficiency.

\Cref{ch:channels-gtm-and-sales} sits upstream of the funnel itself: the go-to-market motion
chosen in \cref{sec:gtm-motion} determines which funnel the firm builds. A product-led growth
(PLG) motion fills the top of the funnel with self-serve trial activations; a sales-led motion
fills it with MQL handoffs. For this company, a hybrid product-led-sales motion is optimal
--- PLG for SMB, SDR-assisted for mid-market PQLs. Sales-force sizing flows through
\cref{eq:sales-productivity}: to add \money{2000000} in ARR with \money{500000} quotas at
\qty{80}{\percent} attainment, five account executives are required, at a blended
sales-assisted CAC of approximately \money{180} --- still well within the LTV:CAC floor of
\cref{eq:ltv-cac}. That headcount and compensation cost feed directly into the operating
model of \cref{ch:operating-model} as line items in the S\&M expense projection.

The three chapters of Part~III (\cref{ch:digital-funnel-and-crm}, \cref{ch:paid-acquisition-platforms},
\cref{ch:attribution-and-measurement}) form a closed loop: the funnel quantifies what is
converted; the platform mechanics determine what it costs; attribution confirms whether the
spend was causal. \Cref{ch:channels-gtm-and-sales} extends the loop one stage earlier, to the
structural choice of \emph{how} to reach the market before any dollar of acquisition spend is
committed.

%% ─────────────────────────────────────────────────────────────
\section{Strategy: Position and Defend}\label{sec:capstone-strategy}
%% ─────────────────────────────────────────────────────────────

Good unit economics in a structurally weak market are temporary. \cref{eq:hhi} in
\cref{ch:competitive-analysis} places the B2B SaaS market at approximately
\num{2260} (percentage-point-squared), just below the highly-concentrated threshold ---
moderately concentrated, with meaningful rivalry and low entry barriers. In Porter's
Five Forces language (\cref{ch:competitive-analysis}): buyer power is high (switching costs
are low without integrations), the threat of entry is high (infrastructure costs have
fallen sharply), and rivalry is intense at a \money{50}/month price point where feature
parity is common. The structural diagnosis is clear: the firm cannot win on price or
feature completeness at its current scale. It must build structural barriers.

The moat taxonomy from \cref{ch:growth} identifies switching costs as the primary lever.
Each API integration added to a customer's workflow increases the migration cost of leaving
by an estimated \money{3000}--\money{6000} --- several years of subscription value at
the current price. A deeply integrated customer faces a genuinely irrational defection
decision even if a nominally superior product arrives. The secondary moat is the integration
marketplace's network effect: \cref{eq:metcalfe} shows that at \(n=100\) vendor integrations,
\(n(n-1)/2 = 4{,}950\) unique interaction pairs exist, versus only 1,225 at \(n=50\). Each
new integration makes every existing integration more valuable, which attracts more vendors,
which attracts more customers --- the demand-side compounding that is absent from a pure
product-quality advantage.

The VRIN test from \cref{ch:advantage-and-disruption} applies: the proprietary usage data
accumulated across the integration marketplace is \emph{valuable} (it enables benchmarking
features), \emph{rare} (competitors cannot observe the same cross-customer data set),
\emph{inimitable} (the data accumulates only through operating at scale over time), and
\emph{non-substitutable} (no public dataset replicates what is created by the network).
A data asset satisfying VRIN supports the sustained ROIC\,>\,WACC spread of \cref{eq:roic}
beyond the horizon where any single feature advantage would have eroded.

\Cref{ch:data-and-ai-strategy} operationalises the VRIN test of \cref{sec:rbv-vrin} for the
data layer specifically. The proprietary multi-tenant usage telemetry (\cref{sec:data-asset})
passes all four VRIN tests: the benchmark feature it enables is valuable, the corpus is rare
and inimitable, and no public dataset substitutes for it. The AI build-vs-buy decision
(\cref{sec:ai-build-buy}) applies \cref{eq:npv} to determine that fine-tuning on this corpus
is only positive-NPV if retention improves by roughly \money{50} per customer, not just a
marginal quality increment. AI governance (\cref{sec:ai-governance}) frames model risk as an
agency problem and concentrates human-in-the-loop review where the product of error probability
and CLV is largest. ESG materiality (\cref{sec:esg-data}) is treated as a measurement problem
with double-materiality boundaries, where the primary financial mechanism at \money{4000000}
ARR is WACC reduction through reduced fundraising friction, not NOPAT directly.

The strategy section's financial output feeds directly into valuation: the moat justifies
a non-zero terminal growth rate $g$ in \cref{eq:terminal-value}. Without a moat,
competition compresses $g$ toward zero; with an integration network effect and a
VRIN-positive data flywheel, $g$ can persist at \qty{3}{\percent} or above through the
terminal period. \textcite{porter1985} grounds the competitive-analysis framework;
\textcite{dranove7theconomics} connects HHI and moat strength to sustainable margins.

%% ─────────────────────────────────────────────────────────────
\section{Finance: Value Quantified}\label{sec:capstone-finance}
%% ─────────────────────────────────────────────────────────────

\cref{ch:risk-and-capital-structure} grounds the discount rate. \cref{eq:capm} with
\(\beta=1.4\), a \qty{4}{\percent} risk-free rate, and a \qty{5}{\percent} equity risk
premium yields a cost of equity of \qty{11}{\percent}. Because the firm is equity-financed,
\cref{eq:wacc} collapses to that same \qty{11}{\percent}. This single number propagates
through every present-value calculation in the book: it is the denominator of \cref{eq:clv},
the discount rate in \cref{eq:dcf}, and the growth-rate benchmark in \cref{eq:terminal-value}.

\cref{eq:dcf} in \cref{ch:valuation} discounts five explicit forecast years of free cash
flow. Year-1 FCF is \money{600000}; the trajectory grows at \qty{30}{\percent} in years 1--3
then decelerates to \qty{10}{\percent} in years 4--5, producing a sum of explicit-year
present values of approximately \money{3380000}. The terminal value is the larger component:

\begin{example}[Terminal-value calculation]
\cref{eq:terminal-value} at year-5 normalised FCF of \money{1263748} (\money{1226940}
grown one period at \qty{3}{\percent}), discounted at \(\mathrm{WACC}-g = 0.11-0.03 = 0.08\):
\[
  \mathrm{TV} \;=\; \frac{\money{1263748}}{0.08} \;\approx\; \money{15800000}.
\]
Adding the PV of the terminal value (\(\approx\money{9400000}\) after discounting 5 years
at \qty{11}{\percent}) to the explicit-year sum gives enterprise value
\(\approx\money{12800000}\). Terminal value is \qty{74}{\percent} of the total --- it is
a bet on the steady-state, not on the next five years.
\end{example}

Scenario analysis via \cref{eq:scenario-ev} converts this point estimate into a probability-weighted
distribution. Three scenarios --- bear (\(p=0.25\), EV \money{8000000}), base (\(p=0.50\),
EV \money{15000000}), bull (\(p=0.25\), EV \money{28000000}) --- yield:
\[
  EV^* \;=\; 0.25\times\money{8000000}
           + 0.50\times\money{15000000}
           + 0.25\times\money{28000000}
         \;=\; \money{16500000}.
\]
The \money{3700000} gap between the single-point DCF (\money{12800000}) and the expected EV
(\money{16500000}) reflects the bull scenario's weight: its upside spread exceeds the bear's
downside, so probability-weighting lifts the expectation above the conservative point estimate.
Capturing the bull case --- platform network effects --- is as valuable as managing the
bear-case risk of moat erosion from data-portability regulation.

\Cref{ch:operating-model} builds the DCF's explicit free cash flows from the ground up rather
than taking them as given. \Cref{eq:cohort-revenue} stacks the survival mechanics of every
monthly cohort into an ARR forecast --- the \money{4000000} ARR itself is the output of 38
months of \num{200} new customers per month at \qty{0.65}{\percent} monthly churn. The
three-driver EBITDA model of \cref{eq:op-model} then converts ARR and cost rates into the FCF
that feeds \cref{eq:dcf}. By Year~3, with S\&M falling from \qty{50}{\percent} to
\qty{35}{\percent} of ARR as the channel matures, the EBITDA margin reaches \qty{+7}{\percent}
and the Rule-of-40 score climbs to 37 --- approaching the benchmark threshold. Sensitivity
analysis via \cref{eq:sensitivity} shows that a \money{100} rise in CAC reduces annual EBITDA
by \money{240000}; a \qty{0.35}{\percent} churn increase reduces the steady-state revenue
ceiling by \money{1400000} --- an order of magnitude larger than a year of CAC overspend.
The operating model converts these sensitivities into the bear, base, and bull FCF trajectories
that \cref{eq:scenario-ev} probability-weights.

\Cref{ch:fundraising-and-dilution} turns the \money{12800000} DCF enterprise value into a
Series~A pre-money. Investors do not use DCF in isolation; they apply \cref{eq:arr-multiple}
to anchor the pre-money to market comparables. At ARR \money{4000000} and a Rule-of-40 score
of 15 (Year~1), the market multiple is discounted from \(8\times\) to \(6\times\), yielding a
market-implied pre-money of \money{24000000} --- comfortably between the base (\money{15000000})
and bull (\money{28000000}) scenarios of \cref{eq:scenario-ev}. The \money{12800000} DCF
point estimate thus serves as a conservative floor; the probability-weighted EV of
\money{16500000} and the \(6\times\) market multiple both sit above it. The SAFE outstanding
--- \money{500000} at a \money{5000000} cap, \qty{20}{\percent} discount (\cref{eq:safe-conversion})
--- converts at the cap, giving the SAFE investor \qty{10}{\percent} pre-Series-A, before the
new institutional money is admitted. Runway of 8 months against a 4--6-month process means the
fundraising trigger has already been crossed: \cref{eq:runway} with \money{2500000} cash and
\money{310000} net burn leaves no buffer for delay.

\textcite{damodaran2012} provides the most systematic treatment of terminal-value sensitivity;
\textcite{brealey2020} covers scenario break-even as a complement to point-estimate DCF.

%% ─────────────────────────────────────────────────────────────
\section{Execution: Constraints and Incentives}\label{sec:capstone-execution}
%% ─────────────────────────────────────────────────────────────

Financial value is a present-value claim on future execution. The operations lens asks
whether the processes exist to deliver it. \cref{eq:littles-law} in \cref{ch:operations}
states \(L = \lambda W\): with 40 open deals in the mid-market pipeline and 10 closed per
month, the average sales cycle is exactly 4 months. Compressing it to 2 months requires
either halving the pipeline (eliminating unqualified deals) or doubling the close rate ---
but not adding SDR headcount to inflate \(L\) while \(\lambda\) is fixed, which would
lengthen \(W\) mechanically.

On the growth-efficiency side, \cref{ch:business-models} provides two real-time tests.
In a given month, the SaaS adds \money{30000} in new MRR and \money{20000} in expansion;
it loses \money{5000} to contraction and \money{22000} to churn. \cref{eq:quick-ratio} gives:
\[
  \text{Quick Ratio} \;=\; \frac{\money{30000}+\money{20000}}{\money{5000}+\money{22000}}
    \;=\; \frac{50000}{27000} \;\approx\; 1.85.
\]
Growth is positive but the denominator (\money{27000} lost) nearly offsets the expansion
inflow (\money{20000}), revealing that the existing base is close to self-defeating. The
primary bottleneck is churn, not new acquisition.

\cref{eq:rule40} confirms the growth-profitability profile: at \qty{30}{\percent} ARR
growth and \(\qty{-15}{\percent}\) EBITDA margin, Rule of 40 scores \(30 + (-15) = 15\)
--- below the \qty{40}{\percent} benchmark but consistent with a capital-efficient
investment phase. \cref{eq:burn-multiple} at \money{1500000} net cash burned against
\money{1200000} net new ARR gives a burn multiple of \(1.25\times\), which is strong:
the business generates more than one dollar of new ARR for every dollar-and-a-quarter
consumed, well below the caution threshold of 2.

The organisational risk is the incentive misalignment that \cref{ch:organizational-behavior}
names after Kerr: rewarding what is measurable (gross new MRR, funnel volume) rather than
what matters (net economic profit, ROIC). A sales team compensated on closed deals has no
incentive to avoid poor-fit customers with high churn probability; a product team rewarded
on feature velocity has no incentive to build the integration depth that reduces churn.
Board governance (\cref{sec:board-governance}) provides the structural backstop: a board
constituted with independent members, an audit committee, and compensation tied to long-run
ROIC --- not to quarterly ARR bookings --- creates the oversight architecture that holds the
Kerr incentive logic together. Aligning the board, management, and individual contributor
incentives with the metric that unifies the whole system --- ROIC, or its per-customer proxy
LTV:CAC --- resolves the misalignment structurally rather than by managerial monitoring.

%% ─────────────────────────────────────────────────────────────
\section{The Founder: Influence, Story \& Team}\label{sec:capstone-founder}
%% ─────────────────────────────────────────────────────────────

Every number above is produced by people the founder must persuade, fund, and lead. Part~VII
treats those capabilities not as soft skills but as inputs to the same ROIC test --- each one
moves a variable the finance chapters hold fixed.

Persuasion and negotiation (\cref{ch:persuasion-influence}) set where, inside the bargaining
range, value is captured. The Series~A term sheet is a negotiation whose zone of possible
agreement (\cref{eq:zopa}) runs from the founder's reservation value (\cref{eq:reservation})
--- the \money{15000000} DCF floor below which raising destroys more ownership than it buys ---
up to the investor's ceiling near the \money{28000000} bull case. Where the round prices inside
that band is pure distributive surplus: a founder who anchors on the \(6\times\)-ARR comparable
rather than the bear case captures millions in retained ownership that no operating improvement
could match. The same chapter's trust equation (\cref{eq:trust}) governs the demand side ---
buyers who arrive pre-trusted through referrals convert below the \money{600} CAC target ---
and is the reason the manipulative ``laws of power'' genre is rejected outright: eroded trust
raises churn and collapses the very CLV the whole model rests on.

Storytelling (\cref{ch:storytelling-communication}) is the delivery mechanism for that
persuasion. The pitch that converts the \money{15000000} floor into a funded \money{24000000}
pre-money is a narrative, not a spreadsheet read aloud; the same narrative discipline lifts
trial-to-paid conversion at the binding funnel gate of \cref{eq:funnel-chain} and seeds the
word-of-mouth that lowers blended CAC. A true value proposition that cannot be transmitted is,
to the market, indistinguishable from one that does not exist.

Founder psychology (\cref{ch:founder-psychology}) is the discipline that holds the plan together
across the eight-month runway, and its claims are the ones that survive meta-analysis:
self-efficacy predicts performance, specific and difficult goals beat vague intention,
task-relevant human capital predicts venture success more than generic experience, and
implementation-intention habits convert resolve into routine. These are precisely what sustain
the execution discipline that \cref{eq:quick-ratio} and \cref{eq:burn-multiple} measure. The
enterprise value of \money{12800000} is built from cohort-survival curves, learning-phase signal
density, and an explicit scenario distribution --- it is realised by a founder who installs those
routines, not one who visualises the outcome.

Leadership (\cref{ch:leadership-charisma}) closes the loop back to ROIC through the workforce.
Transformational leadership and the trainable charismatic tactics that
\cref{ch:storytelling-communication} supplies recruit and retain the five account executives the
operating model of \cref{ch:operating-model} assumes; psychological safety determines whether
the churn-reduction finding of the execution section is implemented or quietly ignored; and the
founder's-dilemma decisions about co-founder equity splits set the cap table that
\cref{eq:safe-conversion} later dilutes. A team that cannot be led at \money{4000000} ARR cannot
deliver the plan the \money{24000000} pre-money is pricing. Leadership is not adjacent to the
financial model --- it is the mechanism that produces the NOPAT.

%% ─────────────────────────────────────────────────────────────
\section{Data, AI \& ESG: The 2026 Additions}\label{sec:capstone-new}
%% ─────────────────────────────────────────────────────────────

The three themes added most prominently in the 2026 edition --- data as a strategic asset
(\cref{sec:data-asset}), AI governance (\cref{sec:ai-governance}), and ESG as a
capital-allocation discipline (\cref{sec:esg-capital}) --- are not separate strategic agendas.
They reduce to the same ROIC\,>\,WACC test that the rest of the book applies to every decision.
A data flywheel is strategically valuable only if it generates product improvements that move
NOPAT upward or allow the firm to defend a price premium that sustains contribution margin; if
it does neither, it is an infrastructure cost dressed as a moat. AI governance is worth its
overhead if and only if the expected value of errors caught exceeds the review cost ---
precisely the NPV calculation the book applies to any investment decision, and exactly the
calculation \cref{sec:ai-governance} sets up with CLV as the value-at-risk denominator. ESG
spend creates economic profit only through two specific channels: reducing WACC by lowering
the risk premium investors assign to regulatory or talent disruption, or improving NOPAT by
reducing churn (talent retention) or raising pricing power (brand trust). At \money{4000000}
ARR the dominant channel is WACC reduction through Series~A fundraising friction --- ESG
data quality is a diligence gate, not an altruistic exercise. All three themes are strategic
only when they move NOPAT, WACC, or invested capital. When they do not, they are costs.

%% ─────────────────────────────────────────────────────────────
\section{The Synthesis}\label{sec:capstone-synthesis}
%% ─────────────────────────────────────────────────────────────

The value-creation spine in \cref{fig:value-spine} is not a metaphor; it is a causal chain.
Start at the left.

\emph{Foundations} establish the test: ROIC must exceed WACC. That test is abstract until
the firm fills in its own numbers. For this company, the verdict is \money{270000} of annual
economic profit --- positive, but not secure. The spread is \qty{9}{\percent} today; the
question is whether it persists.

\emph{Customer economics} give the test its operational content. The mid-market segment was
chosen because its LTV potential (\money{3150}) provides enough headroom over a \money{600}
CAC to sustain a 5.25× ratio --- the per-customer expression of ROIC\,>\,WACC. Brand work
lowers the effective CAC over time; pricing discipline keeps the margin structure that
generates the \money{252}/year of contribution. Without both, CLV shrinks and the ratio
falls toward the 3× floor.

\emph{Acquisition} converts the unit-economics model into real demand. The funnel delivers
200 new customers per month at the target cost, but only because the platform's learning
phase has enough signal density to bid intelligently. The acquisition machine is fragile:
fragment the budget across ten ad sets, or set the CPA target below the learning-phase
minimum, and the whole engine stalls. Attribution discipline --- iROAS rather than last-click
ROAS --- ensures that reported efficiency is real, not borrowed from organic demand.

\emph{Strategy} determines how long the economics last. A product without a moat earns
its ROIC\,>\,WACC spread until a competitor replicates the features; typically that window
is two to four years. The integration switching costs and the nascent network effect of
the vendor marketplace extend the window by raising the cost of defection and by
compounding value with each new node. These are not nice-to-haves; they are the mechanism
by which the \qty{3}{\percent} terminal growth rate in the DCF is defensible rather than
arbitrary. Strategy sets $g$.

\emph{Finance} discounts the whole system at the appropriate rate, translating the
operating story into an enterprise value. The DCF is not a valuation model; it is a
discipline that forces every assumption about growth, margin, and durability into a
single number that can be stress-tested. The scenario analysis reveals that the bear case
--- regulation collapses the switching-cost moat, growth stalls, EV falls to \money{8000000}
--- is a \qty{25}{\percent} probability event. Managing that risk explicitly, rather than
assuming the base case, is the difference between finance as a reporting function and
finance as a strategic tool.

\emph{Execution} is where strategy and finance are either realised or not. The quick ratio
of 1.85 and the burn multiple of 1.25 are not grades; they are coordinates. They locate the
business on the efficiency frontier and reveal the next constraint: churn reduction, not
new acquisition volume, is the lever that moves both metrics toward elite territory. Little's
Law says the 4-month sales cycle can be cut to 2 months, but only by reducing unqualified
pipeline, not by adding salespeople to a pipeline the close team cannot clear faster.
Incentive structures must reinforce these findings, or the organisation will optimise the
metric it is paid on --- which is rarely the metric that matters.

\emph{The founder} is the variable the model cannot hold constant. Persuasion sets where in the
bargaining range the round prices; story decides whether the pitch and the funnel convert;
psychology supplies the discipline --- and the refusal of wishful thinking --- that keeps the
forecast honest; leadership recruits and retains the team that turns the plan into NOPAT. These
are not a softer epilogue to the financial chapters; they are the human inputs to every number
the spine carries.

The full loop closes at ROIC. Position determines whether customers are willing to pay the
price that generates the contribution margin. Customer economics determine whether the
acquisition cost is recoverable within the product's life. Acquisition efficiency determines
whether the CAC is actually \money{600} or quietly \money{900}. Strategic durability
determines whether the margin and growth persist long enough to justify the terminal value
that dominates the DCF. Execution discipline determines whether the plan arrives on time.
Every instrument in the book is pointing at the same reading: \(\mathrm{ROIC}>\mathrm{WACC}\)
--- sustained, not just today, but through the terminal horizon the market is pricing.

\textcite{brealey2020}, \textcite{damodaran2012}, \textcite{gupta2004customer},
\textcite{porter1985}, and \textcite{osterwalder2010} are the five pillars whose frameworks
this chapter weaves together. The integration is the book's contribution; the sources are
theirs.
